% sec-02-methods.tex - Section 2, Methods.  FINAL stage.
% Figures 3, 4 and 5 drawn; Tables 6 and 7 populated.  Senior-author pass: the
% Projects column of Table 6 narrows to its widest cell, and two closing
% subsections state the adoption cost and the provenance argument the rest of
% the paper depends on.
\section{Methods}
\label{sec:methods}

The method is one master prompt, decomposed by the model into sub-prompts that
the model then executes in order, with every generated file committed and pushed
before the next is written. Earlier 2026 papers in this repository reached a
similar effect through a sequence of user prompts, one per stage: the
regulatory adaptations of March 2026, the national platform paper, and the site
documentation package were all built that way, and the author had to be present
at each stage boundary to issue the next instruction \cite{natplatform2026}.
Under the master prompt schedule the author is present once, at the start, and
thereafter watches the branch. The schedule used here has eight stages: five
that produce machine-readable figure specifications, one per diagram platform,
and three that produce the paper itself at increasing quality, draft, full and
final. Figure 3 draws that schedule as a single activity and marks where the
concurrency actually is.

Two properties of the schedule matter more than its length. The first is that
the twenty-five figure slots are fixed before the fork, so five specification
stages can run against one plan without any of them invalidating another; the
second is that the unit of work is one file rather than one stage, so the branch
carries a usable state after every write rather than after every stage. Figure 4
draws one turn at that granularity. Table 6 counts how many separate deposited
projects have used each element of the method, which is the basis on which it is
offered as reliable rather than as novel.

% ---------------------------------------------------------------------------
% FIGURE 3.  plantuml-type, activity with fork and join.  Five concurrent
% branches at a 3.30 cm pitch, then a strictly serial chain at 1.10 cm.
% ---------------------------------------------------------------------------
\begin{appfloat}[!tb]
\begin{appfig}[plantuml-type / activity with fork and join]
\node[umlinit] (i0) at (7.50,0) {};
\node[umlbox,text width=44mm] (a1) at (7.50,-0.85) {Master prompt received,\\one turn, all directives};
\node[umlbox,text width=52mm] (a2) at (7.50,-2.05) {Decompose into eight sub-prompts\\and fix the 25 figure slots};
\node[umlbar,minimum width=13.6cm] (fk) at (7.50,-2.95) {};
\node[umlctrl,text width=25mm] (s1) at (0.90,-4.25) {Stage 1 mermaid\\6 specifications};
\node[umlctrl,text width=25mm] (s2) at (4.20,-4.25) {Stage 2 plantuml\\4 specifications};
\node[umlctrl,text width=25mm] (s3) at (7.50,-4.25) {Stage 3 d2\\6 specifications};
\node[umlctrl,text width=25mm] (s4) at (10.80,-4.25) {Stage 4 diagrams\\4 specifications};
\node[umlctrl,text width=25mm] (s5) at (14.10,-4.25) {Stage 5 graphviz\\5 specifications};
\node[umlbar,minimum width=13.6cm] (jn) at (7.50,-5.55) {};
\node[mmdec,aspect=2.2,text width=24mm] (d1) at (7.50,-6.60) {All 25 slots specified?};
\node[umlbox,text width=40mm] (p6) at (5.10,-7.85) {Stage 6 draft new trial\\instructions, slots, contents};
\node[umlbox,text width=40mm] (p7) at (5.10,-8.95) {Stage 7 full new trial\\every figure drawn,\\every table populated};
\node[umlkey,text width=40mm] (p8) at (5.10,-10.25) {Stage 8 final new trial\\senior author pass,\\no publication directory};
\node[umlbox,text width=32mm] (rt) at (12.30,-7.85) {Return to the\\unfilled stage};
\node[umlfinal] (df) at (12.30,-9.05) {};
\draw[trialchar,line width=0.6pt] ([shift={(-0.13,-0.13)}]df.center) -- ([shift={(0.13,0.13)}]df.center);
\node[umlbox,text width=40mm] (ru) at (5.10,-11.55) {Repository update,\\changelog, releases, version};
\node[umlfinal] (fin) at (5.10,-12.45) {};
\fill[trialburg] (fin.center) circle (1.1mm);
% edges
\draw[umlarrow] (i0) -- (a1);
\draw[umlarrow] (a1) -- (a2);
\draw[umlarrow] (a2) -- (fk);
\foreach \n in {s1,s2,s3,s4,s5}{\draw[umlarrow] (\n |- fk.south) -- (\n.north);
  \draw[umlarrow] (\n.south) -- (\n |- jn.north);}
\draw[umlarrow] (jn) -- (d1);
\draw[umlarrow] (d1.west) -- ++(-1.05,0) -- (5.10,-6.60) -- (p6.north)
  node[umlguard,pos=0.55] {yes};
\draw[umlarrow] (p6) -- (p7);
\draw[umlarrow] (p7) -- (p8);
\draw[umlarrow] (p8) -- (ru);
\draw[umlarrow] (ru) -- (fin);
\draw[umlarrow] (d1.east) -- (12.30,-6.60) -- (rt.north)
  node[umlguard,pos=0.45] {no};
\draw[umldash] (rt) -- (df);
\pnote{-0.30}{-13.35}{The fork is honest: the five specification stages share only the 25-slot
figure plan fixed before it, so none can\\invalidate another. The three paper
stages are drawn serially because each reads the whole of its predecessor.}
\end{appfig}
\vspace{-0.6cm}
\figcaption{Figure 3. The eight-stage build as one activity: five diagram stages
fork and\\join on a shared specification, then three paper stages run strictly in
order.}
\end{appfloat}

\begin{apptable}
\begin{tabularx}{\textwidth}{>{\raggedright\arraybackslash}p{4.65cm} >{\raggedright\arraybackslash}p{1.45cm} >{\raggedright\arraybackslash}p{2.45cm} >{\raggedright\arraybackslash}X}
\headrow \thead{Method element} & \thead{Projects} & \thead{First use} & \thead{Evidence}\\
\midrule
Master prompt, single turn & 3 & Jul 2026 & \appfile{funding/pdac-funding-applications/prompts} \\
AI-generated sub-prompt schedule & 3 & Jul 2026 & \appfile{funding/capitalization-plan/sub-prompts} \\
Five-platform specification set & 3 & Jul 2026 & Five directories per project \\
Draft, full, final staging & 8 & Mar 2026 & Every paper directory since March 2026 \\
One file, one commit, pushed live & 3 & Jul 2026 & The branch history of each build \\
Machine-readable figure reuse & 3 & Jul 2026 & Specifications re-read, not re-derived \\
AI peer review round & 5 & Nov 2025 & \S7, Table 24 \\
DOI deposit per artifact & 30 or more & 2024 & \appfile{new-trial-system/abstracts/README.md} \\
\end{tabularx}
\end{apptable}
\vspace{-0.6cm}
\tabcap{Table 6. How many separate deposited projects have used each element of\\the
method, and where the evidence for each count is held.}

% ---------------------------------------------------------------------------
% FIGURE 4.  mermaid-type, sequenceDiagram.  Four lifelines, twelve rows at a
% 0.62 cm pitch, one loop frame over rows 5 to 7, two self-call arcs at a fixed
% 0.85 cm loop width to the right of the Claude lifeline.
% ---------------------------------------------------------------------------
\begin{appfloat}[!tb]
\begin{appfig}[mermaid-type / sequenceDiagram]
\node[mmactor,text width=25mm] (au) at (0.9,0) {Author, one human};
\node[mmactor,text width=25mm] (cc) at (5.2,0) {Claude Code Opus 5};
\node[mmactor,text width=25mm] (rp) at (9.5,0) {Repository tree};
\node[mmactor,text width=25mm] (gh) at (13.6,0) {GitHub branch};
\foreach \x in {0.9,5.2,9.5,13.6}{\draw[mmlife] (\x,-0.55) -- (\x,-8.95);}
\node[mmact,minimum height=8.30cm] at (0.9,-4.70) {};
\node[mmact,minimum height=7.30cm] at (5.2,-4.60) {};
% loop frame drawn first so rows sit above it
\begin{scope}[on background layer]
\node[mmlane,minimum width=13.9cm,minimum height=2.35cm] (lp) at (7.25,-4.85) {};
\end{scope}
\node[mmlanetitle,anchor=north west,font=\tiny\sffamily\bfseries] at ([shift={(1.5mm,-1.2mm)}]lp.north west) {loop [once per generated file]};
\seqrow{-0.95}{0.9}{5.2}{mmmsg}{1 master prompt, one turn}
\draw[mmmsg] (5.2,-1.55) to[out=20,in=-20,looseness=7] node[right=1mm,font=\tiny\sffamily] {2 decompose into eight sub-prompts} (5.2,-2.05);
\seqrow{-2.65}{5.2}{9.5}{mmmsg}{3 read source archives}
\seqrow{-3.27}{9.5}{5.2}{mmret}{4 extracted sources, nothing invented}
\seqrow{-4.15}{5.2}{9.5}{mmmsg}{5 write one file}
\seqrow{-4.85}{5.2}{13.6}{mmmsg}{6 commit and push immediately}
\seqrow{-5.55}{13.6}{0.9}{mmret}{7 branch state visible in real time}
\seqrow{-6.45}{0.9}{13.6}{mmedged}{8 optional interrupt, any point}
\draw[mmmsg] (5.2,-7.05) to[out=20,in=-20,looseness=7] node[right=1mm,font=\tiny\sffamily] {9 defect pass over the whole stage} (5.2,-7.55);
\seqrow{-8.05}{5.2}{13.6}{mmmsg}{10 correction commit}
\seqrow{-8.55}{5.2}{13.6}{mmmsg}{11 repository update commit}
\seqrow{-8.95}{13.6}{0.9}{mmret}{12 stage complete, bundle attached}
% elapsed braces
\draw[decorate,decoration={brace,amplitude=4pt},trialchar,line width=0.5pt]
  (-0.75,-0.85) -- (-0.75,-3.35) node[midway,left=3pt,font=\tiny\sffamily,align=right] {read\\phase};
\draw[decorate,decoration={brace,amplitude=4pt},trialchar,line width=0.5pt]
  (-0.75,-3.75) -- (-0.75,-6.05) node[midway,left=3pt,font=\tiny\sffamily,align=right] {write\\loop};
\pnote{-0.75}{-9.75}{The loop frame encloses exactly three rows, which is the claim: the unit of
work is one file, not one stage. The\\author interrupt at row 8 is dashed and
runs 0.60 cm below the frame, so it crosses no activation bar.}
\end{appfig}
\vspace{-0.6cm}
\figcaption{Figure 4. One generation turn from master prompt to pushed commit,
with the\\four participants, the two return paths, and the point the author can
stop it.}
\end{appfloat}

The second half of the method is storage. Every figure in this paper exists
first as a specification: a perspective statement, a two-line caption, valid
source in the platform's own syntax, a construction table of absolute
coordinates, an edge-routing paragraph, and a value table naming the repository
file each number came from. That format costs about 6 to 9 kilobytes per figure
and can be read back as the input to the next paper; a raster of the same figure
costs 180 to 900 kilobytes and can be read back as nothing. Figure 5 puts the
two records side by side, and Table 7 states what each stored field buys.

% ---------------------------------------------------------------------------
% FIGURE 5.  graphviz-type, dot record.  Two seven-field records at one 0.42 cm
% field height, two outcome records, three edges, no crossings.
% ---------------------------------------------------------------------------
\begin{appfloat}[!tb]
\begin{appfig}[graphviz-type / dot record]
\def\fh{0.42}
\node[gvcellh,minimum width=5.2cm,minimum height=0.50cm,anchor=north west] at (0,0) {Machine-readable specification};
\foreach \i/\f in {1/{perspective statement},2/{caption, two lines, exact},
  3/{source in a fenced block},4/{TikZ construction table},
  5/{edge routing paragraph},6/{value table with sources},7/{repository source list}}{
  \node[gvcells,minimum width=5.2cm,minimum height=\fh cm,anchor=north west]
    at (0,{-0.50-\fh*(\i-1)}) {\f};}
\node[gvcells,minimum width=5.2cm,minimum height=\fh cm,anchor=north west,font=\tiny\bfseries]
  at (0,{-0.50-\fh*7}) {about 6 to 9 kilobytes};
\node[gvcellh,minimum width=5.2cm,minimum height=0.50cm,anchor=north west] at (0,-4.35) {Raster image of the same figure};
\foreach \i/\f in {1/{pixels},2/{no perspective},3/{no caption},4/{no coordinates},
  5/{no routing},6/{no values},7/{no sources}}{
  \node[gvcellg,minimum width=5.2cm,minimum height=\fh cm,anchor=north west]
    at (0,{-4.85-\fh*(\i-1)}) {\f};}
\node[gvcellg,minimum width=5.2cm,minimum height=\fh cm,anchor=north west,font=\tiny\bfseries]
  at (0,{-4.85-\fh*7}) {about 180 to 900 kilobytes};
\node[gvcellh,minimum width=5.0cm,minimum height=0.50cm,anchor=north west] at (8.05,-0.85) {Readable back as input};
\foreach \i/\f in {1/{regenerate at another size},2/{re-color to a new palette},
  3/{extend by one node},4/{diff against the prior version},
  5/{cite the value's source},6/{audit the routing claim}}{
  \node[gvcells,minimum width=5.0cm,minimum height=\fh cm,anchor=north west]
    at (8.05,{-1.35-\fh*(\i-1)}) {\f};}
\node[gvcellh,minimum width=5.0cm,minimum height=0.50cm,anchor=north west] at (8.05,-5.20) {Not readable back};
\node[gvcellg,minimum width=5.0cm,minimum height=\fh cm,anchor=north west] at (8.05,-5.70) {re-drawn by hand};
\node[gvcellg,minimum width=5.0cm,minimum height=\fh cm,anchor=north west] at (8.05,-6.12) {or re-generated from nothing};
\draw[gvedgeb] (5.2,-1.90) -- node[above,font=\tiny,fill=trialwhite,inner sep=1pt] {six fields survive} (8.05,-2.15);
\draw[gvedgeg] (5.2,-6.05) -- node[above,font=\tiny,fill=trialwhite,inner sep=1pt] {none survive} (8.05,-5.91);
\draw[gvedged] (2.60,-3.90) -- node[right,font=\tiny] {rendered once, if ever} (2.60,-4.35);
\node[gvboxm,text width=40mm,anchor=north west] at (8.05,-7.05)
  {Twenty-five specifications: about 190 kilobytes in total};
\pnote{0}{-8.55}{Both left-hand records carry exactly seven fields at one 0.42 cm field height,
so the comparison is field for\\field and the only variable is what each field
holds. No edge enters a record at a field center.}
\end{appfig}
\vspace{-0.6cm}
\figcaption{Figure 5. Two records for one figure: the machine-readable
specification\\and the raster image, with the six fields only the first returns as
input.}
\end{appfloat}

\begin{apptable}
\begin{tabularx}{\textwidth}{>{\raggedright\arraybackslash}p{3.25cm} >{\raggedright\arraybackslash}p{3.45cm} >{\raggedright\arraybackslash}p{2.95cm} >{\raggedright\arraybackslash}X}
\headrow \thead{Property} & \thead{Specification} & \thead{Raster} & \thead{Consequence for the next paper}\\
\midrule
Size on disk & 6 to 9 kilobytes & 180 to 900 kilobytes & The whole figure set fits in one raster's space \\
Caption & Two lines, exact & None & Captions are diffable across versions \\
Coordinates & Absolute, stated & None & A node can be added without a re-layout \\
Routing & Stated per edge & None & Overlap claims are auditable \\
Values & Sourced per cell & None & No number in a figure is unattributable \\
Reuse & Read back as input & Re-drawn by hand & The next paper starts from this one \\
\end{tabularx}
\end{apptable}
\vspace{-0.6cm}
\tabcap{Table 7. What a machine-readable figure specification stores that a\\raster
does not, and what each stored field buys the next paper.}

Two properties of this method are worth separating from its mechanics, because
they are what make it repeatable by somebody else. The first is that nothing in
the schedule requires the author to be a good writer at speed. The specification
stages produce structured objects, not prose, and the paper stages read those
objects; the quality of the figures is therefore a property of the
specifications rather than of the drafting session that used them. The second is
that the branch is the record. There is no separate build log, no note-taking
file, and no reconstruction after the fact: the order in which the twenty-five
specifications, the three stage bundles and the eleven section files appear on
the branch is the order in which they were written, and it is checkable by
anybody with the repository \cite{repo2026}.

\subsection{What the method costs to adopt}

The method is offered as adoptable, so its cost should be stated. It requires
one person, a repository, a model subscription, and a LaTeX toolchain. It does
not require a team, a project manager, a document management system, or a
version control convention beyond one file per commit. The eight-stage schedule
is a text file; the figure specifications are text files; the paper is a text
file set. Nothing in the method is proprietary to the author, and the two
directories that record it, \appfile{new-trial-system/prompts} and
\appfile{new-trial-system/sub-prompts}, are written to be copied rather than
admired.

Three failure modes are worth naming for anybody adopting it. The first is
fixing the figure plan too late: if the twenty-five slots are not settled before
the specification stages fork, the stages contradict one another and the paper
renumbers, which is the layout-instability defect the specifications are written
against. The second is letting a stage hold work back for a batch commit, which
destroys the property that makes the branch a record. The third is writing prose
before the specifications exist, which produces figures that illustrate
paragraphs rather than paragraphs that read figures; every figure in this paper
was specified before the section that discusses it was written, and the
difference is visible in how often the prose says what a figure shows rather
than restating it.

\subsection{Provenance as an output, not a byproduct}

The word provenance is used deliberately throughout this paper. In the prior
system, provenance is an acknowledgement paragraph: a claim about who did what,
made after the fact, unverifiable in principle. In the new system it is the
commit history, which is produced as a side effect of the production rule rather
than written afterward, and which records the order, the timing and the content
of every write. That distinction is what makes the timing claims of
\S\ref{sec:ind} through \S\ref{sec:funding} checkable rather than assertable,
and it is the property a funder should test first, because it is the cheapest to
test and the one everything else depends on.
